\begin{textsample}{POS Dim 1 – vem\_ed\_27 – Score 11.00 – vem\_ed\_27/t142.txt}  \label{ex:f1_pos_019}
Aos Leitores Osvaldo Succi Jr.
Coordenador PCIs Os Projetos Colaborativos Internacionais (PCI / Cesu) \textbf{representam} uma interpretação das Fatecs do Centro Paula Souza (CPS) de uma abordagem de Intercâmbios Virtuais conhecida internacionalmente como COIL (Collaborative Online Internacional Learning – Aprendizagem Internacional Colaborativa realizada on-line).
Resumidamente, professores de Fatecs e de universidades internacionais \textbf{propõem} aos alunos atividades pedagógicas para \textbf{serem} \textbf{desenvolvidas} em equipes mistas (brasileiros e estrangeiros).
O primeiro PCI / Cesu \textbf{foi} realizado em 2013 na Fatec Americana e, em 2014, venceu o Prêmio Santander Universidades de Melhor Parceria Acadêmica.
Em 2024, ao completar dez anos dessa conquista, o Centro Paula Souza \textbf{é} \textbf{reconhecido} como \textbf{referência} em COIL no mundo.
A Jornada de PCIs 2024, realizada em dezembro, celebra esse prestígio e compartilha boas práticas de professores de Fatecs e Instituições de Ensino Superior (IES) brasileiras, bem como \textbf{reflexões} e pesquisas em Intercâmbios Virtuais e \textbf{políticas} para Internacionalização em Casa.
O evento está detalhado na reportagem principal desta edição.
Entre 2013 e 2024, 12.630 fatecanos participaram de 630 PCIs em 53 Fatecs.
As colaborações \textbf{foram} realizadas com 59 IES de 18 países.
A \textbf{expansão} dos PCIs reflete a consolidação das \textbf{ações} \textbf{desenvolvidas} pela equipe dos PCIs / Cesu na prospecção de colaborações, orientação ao \textbf{desenvolvimento} de projetos, acompanhamento e avaliação dos resultados.
E muito desse sucesso vem do engajamento dos professores das Fatecs e do diferencial dos currículos dos Cursos Superiores de Tecnologia do CPS, com aulas de inglês e espanhol na matriz \textbf{curricular}, \textbf{que} facilitam a comunicação em equipes internacionais.
A proeminência das Fatecs do CPS no universo se reflete na visibilidade internacional alcançada em apresentações dos membros da equipe PCI / Cesu em congressos internacionais, a exemplo de Faubai e IVEC; convites para participar do conselhos executivos de redes internacionais (COIL Connect e Red LatAm COIL) e para ministrar cursos de formação na abordagem COIL para docentes e gestores de universidades.
Em dezembro de 2024, estive no México e no Chile, ministrando palestras e capacitações, resumidas nesta edição.
Boa leitura!

% matched lemmas: ação, curricular, desenvolvimento, desenvolvir, expansão, político, propor, que, reconhecer, referência, reflexão, representar, ser
\end{textsample}
